\section{怎样锤炼语言}

不少初学写作的人，在作文时尽管有了明确的中心思想，很好的材料，合理的结构，但是，写出来的文章却干巴巴的，连自己也觉得“没味儿”，甚至明明是很好的意思，却说得含糊不清，词不达意。其主要原因就是在锤炼语言方面没下过苦功。

作文对语言的要求就是必须具备准确性、鲜明性、生动性。具体些说，就是要求用词贴切，造句恰当，能准确地传情达意；语言简洁明快，态度明朗，不含糊暧昧；行文生动形象，摹貌传神。这“三性”是互相关联、不可分割的，其中准确性是基础。只有在准确的基础上才能谈到鲜明、生动。如果离开准确性而片面追求鲜明性和生动性，往往会华而不实，弄巧成拙。一篇好作文，虽不能说字字珠玑，但至少要做到选词准确，炼句精当，富有感情色彩和表现力。要做到这一点也并非易事，必须在作文时，每用一词，每造一句，都要认真推敲，下一番锤炼的功夫。具体说来，应着重注意以下几点。

\textbf{1. 要善于选用最恰当的词语。}

在作文时，要根据字不离词、词不离句、句不离段的原则，仔细推敲字词的含义和色彩，词与词之间的配搭，稍有不妥，就要选用别的更能传情达意的词语。法国著名作家福楼拜曾经说过：“我们不论描写什么事物，要表现它，唯有一个名词；要赋予它运动，唯有一个动词；要得到它的性质，唯有一个形容词。我们必须继续不断地苦心思索，非发现这个唯一的名词、动词与形容词不可；仅仅发现与这些名词、动词或形容词相类似的词语是不行的；不能因为思索困难而去用类似的词语敷衍了事。”我国唐代的伟大诗人杜甫甚至要求自己“语不惊人死不休”。我们尽管达不到这些语言巨匠的高度，但是这种刻意锤炼词语的精神则必须学习。有了这种精神，我们就会自觉地推敲自己作文的词语，努力去选择那唯一的词语，从而不仅可以避免词语使用不当，感情色彩不对，重复罗嗦或生造词语等毛病，而且会逐渐使语言变得准确、简洁、生动起来，有时还颇能传神呢。

例如，有位同学的习作《课间十分钟》里有这样几句话：

“球场上，同学们玩得正热和。小足球被张超一脚踢到了空中，还没有落到地上，又被王海踢了一脚，猛又升了起来。”

他没有“因为思索困难而去用类似的词语敷衍了事”，而认真推敲起来。最后他把“热和”改成“热火朝天”，删去了“一脚”两字；又把“踢”字改成“补”字，“升”字改成“飞”字。我们看，他这一改，语言就准确、生动多了。因为“热和”仅限于亲热，而表达不出热烈的场面，改用“热火朝天”就准确了；既说“踢到空中”，当然是用脚踢的，而且下文还有个“一脚”，为求简炼，删去“一脚”两字很必要；至于把“踢”字改成“补”字，不仅避免了用词前后重复，而且把踢足球的过程和时机表现了出来；特别是最后那个“升”字被改成了“飞”字，形象生动，颇能传神，真是改得好。这说明同学们只要细致观察，用心琢磨，是可以选准那唯一的词语的。

精选词语，正如福楼拜所说，主要是精选那唯一的名词、动词和形容词，因为这些实词最能传情达意。不过，对于那些虽无实意却起语法作用的虚词，我们也不能置之不问。例如一位同学在习作《说“敷衍”》中，先写道：“敷衍，就是待人和做事不认真负责，只作表面上的应付。”再一细想，觉得不妥：待人不恳切，做事不认真，只要沾上其中一项，就是敷衍。因此，他将并列连词“和”改用选择连词“或”。这样一改，判断就周密，表意就准确了。

大诗人杜甫深有体会地说过：“新诗改罢自长吟。”我们写文章也应该这样：每用一词，看看这个词用得恰不恰当，是不是还有表意更准确、更鲜明、更生动的词可以替换。它，改了之后，反复读上几遍，想想是否改得更有味儿，更传神些；至于究竟用哪个词为好，还得要联系上下文进行推敲，直至确定某个词最能传情达意才罢手。

\textbf{2. 要善于把句子写得通顺、简洁。}

句子通顺是锤炼语言的起码条件。句子不通，意思也就无法表达准确，更谈不上鲜明、生动了。值得注意的是，有的同学刻意追求文采，觉得某种说法漂亮，但对这个句子的语法没有把握，仍然勉强去写，结果往往弄出了病句，这真叫弄巧成拙了。

例如有位同学在《雾晨》这篇习作中，这样描写大马路上的雾景：

“纵目望去，自行车群像是来往穿梭的流水，那些公共汽车、面包车、吉普车就是浮沉在河面上的大小船只，而河两岸高大的梧桐树仿佛变成了大雾弥天的神话世界。”

这位同学如果仔细琢磨一下，就会发现：既然是“大雾弥天”，怎能“纵目望”得清楚？谁见过“流水”可以“来往穿梭”？倘若“自行车群”象流水，那些汽车又怎能在这样的“河面”上“浮沉”？“浮沉”是说在水波上忽上忽下，难道那些车辆是上下动荡地前进吗？“梧桐树”又岂能“变成了”“神话世界”？这些不合逻辑或违反语法的句子，表面上似乎很精彩，实际上是站不住脚的。遇到这类情况，宁可写得质朴些，也不要使句子不通。象这段话，我们不妨这样表达：

“透过雾幔，隐约看见各种车辆象两条对流的小河在眼前缓缓而过，汽车的喇叭声、自行车的铃铛声响成一片。路两旁高大的梧桐树则越来越淡地隐没在雾海中。”

这样一换，句子也许平淡些，但表意准确无误；文风朴实，这应是初学写作者首先追求的东西。

句子通顺之外，还得简洁。我们要象鲁迅先生要求的那样，在作文时，竭力将可有可无的字、词、句、段毫不可惜地删去。当用一句话说得清楚的，不要写成两句；当用一个词说得清楚的，不要用两个。这样，语言就干净俐落；否则，就重复罗嗦。

例如，有位同学在写《看电影〈少年犯〉有感》的开头说：“我每当看了一次好电影之后，总是热血沸腾，心潮起伏，总不免要动手动笔，抒写一些自己的心得和感受。”这句话用了四十多个字，其实，只需要“每看一部好电影，我总要写心得”十三字就够了。写“心得”，只能在“看了（或者‘读了’）······之后”，也必定“动手动笔”；“心得”，当然是“自己的”，至于“热血沸腾，心潮起伏”只是套话，用在这儿并不妥贴；而“心得”和“感受”在这句话中是同样的意思，留下一个就行。从全篇作文看，如果直接以“我看了电影《少年犯》，眼前老是浮现着那些少年犯的形象，心头老是想着他们犯罪的原因”开头，那么，这整个一段话都是多余的，可以全部删去。

在追求简洁这个问题上，革命导师马克思堪称是我们学习的典范。据威廉·李卜克内西的回忆，说马克思“对于语言的简洁和正确是一丝不苟的”，“有时到了咬文嚼字的程度。”列宁、斯大林、毛主席的著作，都不仅闪耀着真理的光辉，具有雄辩的逻辑力量，而且在语句的锤炼上无不具有一语破的、不可更易的特色。我国古代的诗人有“为求一字稳，耐得半宵寒”、“二句三年得，一吟双泪流”的名言佳话传世，至今对我们也还是有一定的借鉴作用吧。

当然，文句简洁并非是说愈简愈好，这里有个限度，就是表意明确，没有歧义，又不晦涩难懂。如有个同学写《妹妹的笑》，有这样的句子：“虹美，爱笑，象花。妈妈疼，假吵，也裂嘴笑。”你能懂得他是在说“虹妹长得好看，喜欢笑，一笑起来象花一样。妈妈疼爱她，有时假装不许她笑，可给她一笑也逗得合不拢嘴地笑了”吗？应该明确，文字的繁简是以内容的表达为准则的，而不能为“简”而“简”，更不能不该简的也简。象“也裂嘴笑”究竟是“妈妈”笑呢，还是虹妹笑？“虹美”究竟是名字呢，还是“虹”生得“美”呢？“妈妈疼”什么呢，还是“妈妈”真感到“疼”呢？这样的简法，闹得意思含混不清，就不是简洁而是“简陋”了。有的同学学了些文言文，也学起古人的腔调，写些半文半白的话，语句晦涩古奥，搞得人读不下去，当然也是要不得的。其实，为了表达特定的内容，形成一种特殊的效果，有时还非多用一些词语不可。象鲁迅的《秋夜》的开头写道：“墙外有两棵树，一棵是枣树，还有一棵也是枣树。”读来并不感到繁冗，反而觉得别具意味，能引起人们的注意和会味。如果改成“墙外有两棵枣树”，简倒简了，但也索然无味了。

\textbf{3. 要灵活运用修辞手法，句式富于变化。}

作文的语言能达到用词准确、炼句简洁，已经很好了；如果还能根据文体的特点和表情达意的需要，采用各种富有表现力的修辞手法，选用形象的词语和适当的句式，使文章绘声绘色，摇曳多姿，当然就更好了。

一般说来，由于文体不同，其语言特点也就不完全一致。议论文的语言，侧重要严谨、简要；记叙文的语言，侧重要形象、生动；说明文的语言，侧重要平实、浅显（应用文的语言与说明文的语言大致相同）。其中以记叙文的语言对形象化要求最强烈，最好能达到使读者仿佛觉得作者所写事物的形状和情景就在眼前。要达到这一点，除了观察仔细、想象丰富之外，主要依靠各种修辞手法的巧妙运用。象大家熟悉的比喻、借代、夸张、对偶、排比、拟人等修辞手法，只要用得恰当，就能使语言形象化起来，富有艺术感染力。这里的关键是“恰当”二字。如果使用不恰当，象前面所引的习作《雾晨》那样，连准确性都没有达到，更谈不上什么鲜明生动，提高表情达意的效果了。

另外，善于选择形象的词语和适当的句式，也能使文章显得色彩鲜明，感情洋溢，可以表达出不同的情调、语气和画面。但是要注意：决不能误认为语言要形象，文章要生动，就得多多地堆砌形容词，或者多多使用不同的句式。选择什么样词语或必要的句式，都是由表达什么样的思想内容决定的。许多优秀的文章倒是很少使用形容词，它只是选择富于形象化的词语和恰当的句式，把事物的特征不加渲染烘托地表现出来，就很形象，就很传神。这正是鲁迅先生一向推崇的“白描”手法。例如课文《王冕》中开头的一段，这样来写雨后景色：

“须臾，浓云密布，一阵大雨过了，那黑云边上镶着白云，渐渐散去，透出一派日光来，照耀得满湖通红。湖边山上，青一块，紫一块，绿一块。树枝上都象水洗过一番的，尤其绿得可爱。湖里十来枝荷花，苞子上清水滴滴，荷叶上水珠滚来滚去······”

这里只抓住当时当地雨后景色中最有特征的几点，把云彩、日光、水色、山景、树枝、荷花的形象和色彩，描绘得象画面一样摆在读者面前，语言洗练极了，形象极了，但是它并没有堆砌什么形容词。句子大多简短，间或插上较长的句子；而写雨景、云散、日出、山色，句式参差，节奏明快；至于写树枝、荷花、花苞、荷叶，句式多变，各不相同。作者就这样似乎毫不费力地勾勒出一幅笼罩在山光水色中的雨后翠荷图，传达出热爱自然美景的陶醉情怀。这种写法，很值得我们效法。

议论文的语言，虽然不追求形象的描绘，但是论述仍然要求生动。这就需要善于采用各种恰当的修辞手法，善于灵活运用谚语、成语、典故、古诗词、名言警句等，以增强议论和说理的效果；在论述过程中，还要努力提炼出含蓄、精辟的句子，使文章言简意赅，一语中的。例如，马克思说：“在科学上没有平坦的大道，只有不畏劳苦沿着陡峭山路攀登的人，才有希望达到光辉的顶点。”这个句子，以攀登高峰作比喻，生动地说明在科学方面要取得光辉成就，就要付出艰苦的劳动。句子先用“没有”，再用“只有”、“才有”进行对照，很肯定，很形象而又很严密地说明了这个道理，比起抽象地论述“科学成就来自艰苦的劳动”，给读者的印象不知深刻多少倍。正是由于这句话用了形象的词语和恰当的句式，才使所表达的内容显得如此精警，如今已成了人们经常引用的格言。

一般说明文要求语言平实确切，不追求文字的感情色彩，不用夸饰的语言；为了使读者容易理解所要说明的事物，文字必须写得深入浅出，通俗明白。不过，在特殊类型的说明文中，例如科学小品，对语言的要求则有所不同，因为科学小品是一种把知识性与趣味性融于一体的说明文，因此它对事物的说明采用文艺笔调，如课文中的《琥珀》、《蝉》等，语言就比一般说明文要生动形象得多。总之，我们写作时，应该使语言准确、鲜明、生动，这是锤炼语言的努力方向。要达到这一点，除了上面所说的以外，还得平时多积累丰富的词汇，多作变化句型的练习，特别是有意识地学习活在人民口头中的语言，这样才可能在行文时得心应手，写出语言精彩的文章来。
\newpage